\textbf{Phần 1. Trắc nghiệm lựa chọn}
\setcounterpageref{ex}{0}

\begin{ex} %[2D12V61]
Kết quả của phép tính $256 + 178$ là
\choice
{\True $434$}
{$424$}
{$444$}
{$334$}
\loigiai{
$256 + 178 = 434$.
\textbf{Đáp án của câu hỏi là A}
}
\end{ex}

\begin{ex} %[2D12V61]
Tính tổng của $467$ và $285$.
\choice
{$742$}
{$762$}
{\True $752$}
{$652$}
\loigiai{
$467 + 285 = 752$.
\textbf{Đáp án của câu hỏi là C}
}
\end{ex}

\begin{ex} %[2D12V61]
Điền chữ số thích hợp vào dấu $*$ để: $3*8 + 245 = 613$.
\choice
{$5$}
{\True $6$}
{$7$}
{$4$}
\loigiai{
$613 - 245 = 368$, nên $* = 6$.
\textbf{Đáp án của câu hỏi là B}
}
\end{ex}

\begin{ex} %[2D12V61]
Một trại chăn nuôi xuất bán $375$ con gà trong tuần thứ nhất và $486$ con gà trong tuần thứ hai. Cả hai tuần trại xuất bán bao nhiêu con gà?
\choice
{$851$ con}
{$871$ con}
{$761$ con}
{\True $861$ con}
\loigiai{
$375 + 486 = 861$ (con gà).
\textbf{Đáp án của câu hỏi là D}
}
\end{ex}

\begin{ex} %[2D12V61]
Tìm $x$ biết: $x - 368 = 257$.
\choice
{\True $625$}
{$615$}
{$635$}
{$525$}
\loigiai{
$x = 257 + 368 = 625$.
\textbf{Đáp án của câu hỏi là A}
}
\end{ex}

\begin{ex} %[2D12V61]
Đoạn đường từ tỉnh A đến tỉnh B dài $285\text{ km}$, từ tỉnh B đến tỉnh C dài $348\text{ km}$. Quãng đường từ A qua B đến C dài bao nhiêu ki-lô-mét?
\choice
{$623\text{ km}$}
{$643\text{ km}$}
{\True $633\text{ km}$}
{$533\text{ km}$}
\loigiai{
$285 + 348 = 633$ (km).
\textbf{Đáp án của câu hỏi là C}
}
\end{ex}

\begin{ex} %[2D12V61]
Tính giá trị của biểu thức: $189 + 245 + 316$.
\choice
{$740$}
{\True $750$}
{$760$}
{$650$}
\loigiai{
$189 + 245 + 316 = 750$.
\textbf{Đáp án của câu hỏi là B}
}
\end{ex}

\begin{ex} %[2D12V61]
Số hạng thứ nhất là $458$, số hạng thứ hai lớn hơn số hạng thứ nhất $75$ đơn vị. Tổng của hai số đó là
\choice
{$981$}
{$533$}
{$1 001$}
{\True $991$}
\loigiai{
Số hạng thứ hai: $458 + 75 = 533$. Tổng: $458 + 533 = 991$.
\textbf{Đáp án của câu hỏi là D}
}
\end{ex}

\begin{ex} %[2D12V61]
Điền dấu thích hợp: $367 + 258 \dots 195 + 430$.
\choice
{\True $=$}
{$>$}
{$<$}
{Không so sánh được}
\loigiai{
$367 + 258 = 625$; $195 + 430 = 625$. Vậy hai biểu thức bằng nhau.
\textbf{Đáp án của câu hỏi là A}
}
\end{ex}

\begin{ex} %[2D12V61]
Một cửa hàng buổi sáng bán được $275\text{ kg}$ ngô, buổi chiều bán được $368\text{ kg}$ ngô. Cả ngày cửa hàng bán được bao nhiêu ki-lô-gam ngô?
\choice
{$633\text{ kg}$}
{$653\text{ kg}$}
{\True $643\text{ kg}$}
{$543\text{ kg}$}
\loigiai{
$275 + 368 = 643$ (kg).
\textbf{Đáp án của câu hỏi là C}
}
\end{ex}

\begin{ex} %[2D12V61]
Tổng của số lẻ bé nhất có ba chữ số khác nhau và $589$ là
\choice
{$682$}
{\True $692$}
{$702$}
{$690$}
\loigiai{
Số lẻ bé nhất có 3 chữ số khác nhau là 103. Tổng: $103 + 589 = 692$.
\textbf{Đáp án của câu hỏi là B}
}
\end{ex}

\begin{ex} %[2D12V61]
Bác Nam có hai thùng mật ong. Thùng thứ nhất chứa $185\text{ l}$, thùng thứ hai chứa $235\text{ l}$. Bác đã bán đi $120\text{ l}$. Bác còn lại bao nhiêu lít mật ong?
\choice
{$290\text{ l}$}
{$310\text{ l}$}
{$420\text{ l}$}
{\True $300\text{ l}$}
\loigiai{
Tổng số mật ong: $185 + 235 = 420$ l. Còn lại: $420 - 120 = 300$ l.
\textbf{Đáp án của câu hỏi là D}
}
\end{ex}

\textbf{Phần 2. Trắc nghiệm đúng sai}
\setcounterpageref{ex}{0}

\begin{ex} %[2D12V61]
Xét các phép tính cộng có nhớ trong phạm vi 1 000:
\choiceTF
{\True 378 + 245 = 623.}
{489 + 12 = 591.}
{\True Khi cộng có nhớ sang hàng chục, ta phải thêm 1 vào hàng chục của tổng.}
{\True Tổng của 450 và 550 bằng đúng 1 000.}
\loigiai{
\begin{itemize}
    \item Ý a) Đúng: 8 + 5 = 13 viết 3 nhớ 1; 7 + 4 = 11 nhớ 1 là 12 viết 2 nhớ 1; 3 + 2 = 5 nhớ 1 là 6. Đúng.
    \item Ý b) Sai: 489 + 12 = 501, không phải 591.
    \item Ý c) Đúng: Quy tắc cộng có nhớ. Đúng.
    \item Ý d) Đúng: 450 + 550 = 1 000. Đúng.
\end{itemize}
\begin{itemize}
    \item a) \textbf{Đúng}.
    \item b) \textbf{Sai}.
    \item c) \textbf{Đúng}.
    \item d) \textbf{Đúng}.
\end{itemize}
}
\end{ex}

\begin{ex} %[2D12V61]
Trường Tiểu học Hòa Bình có 368 học sinh nam và 357 học sinh nữ:
\choiceTF
{\True Trường có tất cả 725 học sinh.}
{\True Số học sinh nam nhiều hơn số học sinh nữ 11 bạn.}
{Nếu có thêm 5 bạn nữ chuyển đến thì số nam bằng số nữ.}
{\True Tổng số học sinh ít hơn 800 bạn.}
\loigiai{
\begin{itemize}
    \item Ý a) Đúng: 368 + 357 = 725 học sinh. Đúng.
    \item Ý b) Đúng: 368 - 357 = 11 bạn. Đúng.
    \item Ý c) Sai: 357 + 5 = 362 < 368 (vẫn ít hơn). Sai.
    \item Ý d) Đúng: 725 < 800. Đúng.
\end{itemize}
\begin{itemize}
    \item a) \textbf{Đúng}.
    \item b) \textbf{Đúng}.
    \item c) \textbf{Sai}.
    \item d) \textbf{Đúng}.
\end{itemize}
}
\end{ex}

\textbf{Phần 3. Trả lời ngắn}
\setcounterpageref{ex}{0}

\begin{ex} %[2D12V61]
Tính tổng: $476 + 389$.
\shortans{865}
\loigiai{
\textbf{Đáp số:} $865$\\
$476 + 389 = 865$.
}
\end{ex}

\begin{ex} %[2D12V61]
Tìm $x$ biết $x - 287 = 436$.
\shortans{723}
\loigiai{
\textbf{Đáp số:} $723$\\
$x = 436 + 287 = 723$.
}
\end{ex}

\begin{ex} %[2D12V61]
Đội một trồng được 345 cây, đội hai trồng được 385 cây. Cả hai đội trồng được bao nhiêu cây?
\shortans{730}
\loigiai{
\textbf{Đáp số:} $730$\\
$345 + 385 = 730$ cây.
}
\end{ex}

\begin{ex} %[2D12V61]
Tổng của hai số là 850. Nếu tăng số hạng thứ nhất 35 đơn vị và tăng số hạng thứ hai 15 đơn vị thì tổng mới là bao nhiêu?
\shortans{900}
\loigiai{
\textbf{Đáp số:} $900$\\
$850 + 35 + 15 = 900$.
}
\end{ex}

\begin{ex} %[2D12V61]
Tính: $195 + 285 + 320$.
\shortans{800}
\loigiai{
\textbf{Đáp số:} $800$\\
$195 + 285 + 320 = 800$.
}
\end{ex}

\begin{ex} %[2D12V61]
Một cuộn dây dài 258 m, cuộn thứ hai dài 367 m. Cả hai cuộn dài bao nhiêu mét?
\shortans{625}
\loigiai{
\textbf{Đáp số:} $625$\\
$258 + 367 = 625$ m.
}
\end{ex}
